
Marine wildlife conservation is an important issue that requires more attention due to the significant decline in the population \cite{nelms2021marine}. One of the key challenges in marine conservation is the accurate identification of marine species, which is often time-consuming and error-prone. With the rise of technology and the increasing availability of marine wildlife imagery, computer vision techniques, such as object detection and classification with deep learning, have emerged as a potential solution for accurate and efficient species identification and monitoring.

However, machine learning techniques are highly dependent on the amount and quantity of the data used. In the case of marine wildlife imagery, obtaining high-quality data is often difficult due to the challenging conditions and often requires a lot of manual works. An imbalanced dataset may occur when the number of images in a particular species is limited compared to others. Additionally, the complexity of marine wildlife aerial images also lies in high resolution and various types of distortions such as perspective distortion and non-uniform illumination. 

To address these challenges, data augmentation techniques could be used to artificially increase the size of a data-set by creating variations of existing images. The primary goal of image augmentation is to increase the diversity of the training data and prevent over-fitting. Finding the optimal augmentation policy is challenging because it heavily depends on the dataset, and there is no heuristic to determine it. 

Therefore, this thesis aims to investigate the effectiveness of various types of augmentations, including geometric, color-space, cutout, and mixed, for marine wildlife imagery in both classification and detection tasks. Additionally, the thesis will investigate the potential use of Generative Adversarial Networks (GANs) to generate images and enhance the classification of underrepresented species.



